\chapter{Das diskrete Logarithmusproblem}

Alle bis heute bekannten Algorithmen liefern keine effizienten
Berechnungsverfahren zur Bestimmung des diskreten Logarithmus eines
Gruppenelements aus einer endlichen zyklischen Gruppe. Damit wird das 
diskrete Logarithmusproblem neben dem Faktorisieren als eines der beiden 
bekannten harten Probleme angesehen, die sich f"ur kryptographische 
Zwecke eignen. Das Verfahren von ElGamal st"utzt sich auf diese Tatsache.

\section{Einf"uhrung}

Es sei $G$ eine endliche zyklische Gruppe und $\alpha$ ein Erzeuger von $G$.
Dann gilt 
$$ G = \left\{\alpha^i\mid 0\le i < \#G\right\}, $$
wobei $\#G$ die Gruppenordnung von $G$ ist. Der diskrete Logarithmus
eines Elements $\beta$ zu der Basis $\alpha$ in $G$ ist eine ganze Zahl $x$ mit
$\alpha^x = \beta$. Wenn $x$ auf das Intervall $0\le x < \#G$ beschr"ankt ist,
dann ist der diskrete Logarithmus von $\beta$ zu der Basis $\alpha$
eindeutig. Wir schreiben $x = \dlog{\alpha}{\beta}$.

Das diskrete Logarithmusproblem besteht darin, einen effizienten
Algorithmus zu finden, der Logarithmen in einer gegebenen Gruppe 
berechnet.

Man kann sofort zwei Algorithmen zur Berechnung der diskreten Logarithmen
in einer endlichen zyklischen Gruppe $G$ angeben. Der erste berechnet
einmal eine Liste der Logarithmen aller Gruppenelemente. Der zweite
berechnet so lange Potenzen von $\alpha$, bis eine "Ubereinstimmung mit
$\beta$ auftritt. Diese trivialen Algorithmen sind wertlos, wenn $\#G$ 
gro"s ist.

\section{Baby-Step Giant-Step Algorithmus}
Der Baby-Step Giant-Algorithmus eignet sich, um Logarithmen in beliebigen
zyklischen Gruppen zu berechnen. Er stellt eine erhebliche Verbesserung
gegen"uber den in der Einf"uhrung beschrieben trivialen Algorithmen, ist
aber praktisch nicht einsetztbar, wenn die Gruppenordnung zu gro"s ist.

Wir wollen den diskreten Logarithmus $\dlog{\alpha}{\beta}$ berechnen.
Es sei $m=\oben{\sqrt{\#G}}$.
\begin{enumerate}
  \item Stelle eine Liste der Paare $(i,\alpha^i)$ f"ur $0\le i < m$
    (offensichtlich gilt $i=\dlog{\alpha}{\alpha^i}$) und sortiere
    diese Liste nach der zweiten Komponente.
  \item Berechne f"ur jedes $j$ mit $0 \le j < m$ $\beta \alpha^{-j m}$ und
    "uberpr"ufe (mit Bin"arsuche), ob dieses Element als zweite 
    Komponente eines Listeneintrags vorkommt.
  \item Falls $\beta\alpha^{-jm} = \alpha^i$ f"ur ein $i$ mit $0 \le i < m$,
    dann gilt $\beta = \alpha^{i+jm}$ und $\dlog{\alpha}{\beta} = i + j m$. 
\end{enumerate}
Der Algorithmus stellt eine Liste mit $\Oklasse{m}$ Eintr"agen und 
ben"otigt $\Oklasse{m \log m}$ Operationen, um die Liste zu sortieren
und dann in der Liste zu suchen. Mit einer Operation ist
eine Gruppenoperation oder ein Vergleich gemeint. Der Name beschreibt in
kennzeichnender Weise die Funktionsweise des Algorithmus.

"Ubung: Wie sieht die Giant-Step Baby-Step Version des Algorithmus?

\section{Index Calculus}

Wir beschreiben generisch die Funktionsweise der Index Calculus Methode.
Es sei $G$ eine endliche zyklische Gruppe der Ordnung $n$, die durch
$\alpha$ erzeugt wird. Wir wollen die Logarithmen der Gruppenelemente zu
dieser Basis berechnen. Es sei $S=\{p_1, p_2, \ldots, p_t\}$ eine
Untermenge von $G$ mit der Eigenschaft, da"s ein "`bedeutender"' Teile
der Gruppenelemente sich als Produkt der Elemente aus $S$ schreiben l"a"st.
Die Menge $S$ wird als die {\em Faktorbasis} bezeichnet.

Im ersten Schritt versuchen wir die Logarithmen aller Elemente aus $S$
zu bestimmen. Dazu w"ahlen wir eine Zufallszahl $a$ und versuchen $\alpha^a$
als Produkt der Elemente aus $S$ zu schreiben:
\begin{equation}
  \alpha^a = \prod_{i=1}^t p_i^{\lambda_i}.\label{produkt}
\end{equation}
Wenn wir so eine Darstellung finden k"onnen, erhalten wir aus \ref{produkt}
eine lineare Kongurenz
\begin{equation}
  a \equiv \sum_{i=1}^t \lambda_i \dlog{\alpha}{p_i} \pmod{n}.\label{relation}
\end{equation}
Nachdem wir eine gen"ugend gro"se Anzahl (gr"o"ser als $t$) an Relationen
der Form~\ref{relation} aufgestellt haben, k"onnen wir erwarten, da"s das 
zugeh"orige lineare
Gleichungssystem eine eindeutige L"osung f"ur die Unbekannten 
$\dlog{\alpha}{p_i}$ mit $1 \le i \le t$ besitzt.

Im zweiten Schritt berechnen wir individuelle Logarithmen in $G$.
Es sei ein $\beta\in G$ gegeben, von dem wir den Logarithmus 
$x=\dlog{\alpha}{\beta}$ bestimmen m"ochten. Wir w"ahlen solange Zufallszahlen
$s$, bis $\alpha^s\beta$ sich als Produkt von Elementen aus $S$ schreiben
l"a"st:
\begin{equation}
  \alpha^s\beta = \prod_{i=1}^t p_i^{b_i}. \label{inS}
\end{equation}
Es gilt dann
$$ \dlog{\alpha}{\beta} \equiv \sum_{i=1}^t b_i\dlog{\alpha}{p_i} - s
   \pmod{n}. $$
Um die Beschreibung der Index Calculus Methode abzuschlie"sen, m"ussen wir
angeben, wie ein geeignetes $S$ zu w"ahlen ist und wie effizient die
Relationen \ref{produkt} und \ref{inS} zu erzeugen sind. Ein geeignetes
$S$ ist klein (das Gleichungssystem im ersten Schritt ist dann nicht zu
gro"s) und gleichzeitig mu"s die Anzahl der Gruppenelemente, die "uber
$S$ faktorisieren, gro"s sein (die erwartete Zeit, um die Relationen
\ref{produkt} und \ref{inS} aufzustellen, ist nicht zu lang). Solche Spezifikationen
sind f"ur einige Endliche K"orper bekannt. Wir beschreiben kurz die
Index Calcus Methode f"ur endliche K"orper. Die Algorithmen machen Gebrauch 
von den unterschiedlichen Darstellungsm"oglichkeiten der K"orperelemente. 
Deswegen wird eine kurze Einf"uhrung in die Theorie der endlichen K"orper
gegeben.

Die erwarteten Laufzeiten der Algorithmen werden durch die Funktion
$$ L[x,\alpha,c] = \Oklasse{\exp((c+o(1))(\log x)^\alpha)
   (\log\log x)^{1-\alpha})}, $$
beschrieben, wobei $x$ die Gruppenordung ist, $0\le\alpha\le 1$ und $c$ Konstanten sind.
Wenn $\alpha = 0$, dann ist $L[x,\alpha,c]$ polynomial in
$\log x$. Bei $\alpha = 1$ ist $L[x,\alpha,c]$ dagegen exponentiell in
$\log x$. Wenn $0 < \alpha < c$, dann ist $L[x,\alpha,c]$ 
{\em subexponentiell} in $\log x$.

\subsection{Endliche K"orper}

  \begin{satz}[Primk"orper]
    F"ur $m\in\N$ sind 
    (i)  $m$ ist eine Primzahl und (ii) $\Z/m\Z$ ist ein K"orper 
    "aquivalent.
  \end{satz}
  Die endlichen K"orper $\koerper{p} = \Z/p\Z$, wobei $p$ eine Primzahl ist,
  hei\3en {\em Primk"orper}.

  \begin{definition}[K"orpererweiterung]
    Unter einer K"orpererweiterung versteht man ein Paar von K"orpern 
    $K \subset L$, wobei $K$ ein Teilk"orper von $L$ ist. Bezeichnung: $L/K$
    oder $K<L$. L hei\3t Erweiterungsk"orper von K.
  \end{definition}
  Es sei $L$ eine K"orpererweiterung von $K$. Man kann die Multipikation
  auf $L$ einschr"anken zu einer Multipikation $K \times L \rightarrow L$  
  und auf diese Weise $L$ als $K$-Vektorraum auffassen. Mit Hilfe dieser
  Beobachtung erh"alt man eine wichtige Strukturaussage "uber endliche
  K"orper. 

  Es sei $\koerper{q}$ ein endlicher K"orper mit q Elementen und 
  $\koerper{q}/\koerper{p}$. Dann ist $\koerper{q}$ ein endlicher 
  $\koerper{p}$-Vektorraum, da $\koerper{q}$ endlich ist. Es sei
  $ \mbox{dim}_{\koerper{p}}(\koerper{q}) = [\koerper{q} : \koerper{p}] = n$
  die Dimension und $\{\omega_1, \ldots, \omega_n\}$ eine Basis von
  $\koerper{q}$ "uber $\koerper{p}$. Jedes Element von $\koerper{q}$ l"a\3t
  sich in eindeutiger Weise als Linerakombination
  $ a_1 \omega_1 + \ldots + a_n \omega_n $
  mit $a_1, \ldots, a_n\in\koerper{p}$ schreiben. Die Anzahl der Elements
  des K"orpers $\koerper{q}$ ist also $p^n$.

  Durch die Vektorrauminterpretation erhalten wir die Strukturaussage, da\3
  die Anzahl der Elemente eines endlichen K"orpers eine Primzahlpotenz
  ist.
  \begin{definition}[Irreduzibel]
    Ein Polynom $f(x)\in K[x], f\neq 0$, wobei $K$ ein K"orper ist, hei\3t 
    {\em irreduzibel}, falls f"ur jede Zerlegung $f(x)=g(x) h(x)$ mit 
    $g(x),h(x)\in K[x]$ 
    gilt $g(x)\in K^*$ oder $h(x)\in K^*$. 
  \end{definition} 

  \begin{satz}[Konstruktion eines K"orpers mit $p^n$ Elementen]
    Es sei $f(x)\in\koerper{p}[x]$ ein irreduzibles Polynom mit $\grad f=n$.
    Dann ist der Restklassenring $\koerper{p}[x]/f(x)$ sogar ein K"orper.
  \end{satz}
  \beweis Die K"orpereigenschaft kann leicht dadurch verdeutlicht werden, 
  indem der erweiterte ggT-Algorithmus ELBA betrachtet wird. Der ELBA
  stellt den ggT zweier Elemente eines euklidischen Rings als eine
  Linearkombination von ihnen dar. 
  Wir fordern vor"ubergehend, da\3 $f(x)$ normiert ist, d.h. der 
  Leitkoeffizient ist gleich $1$. Da $f(x)$ ein normiertes irreduzibles 
  Polynom vom Grad $n$ ist, ist der $\ggt (f(x),g(x))=1$ f"ur alle 
  Polynome $g(x)\in\koerper{p}[x]$ mit $\grad(g(x)) < n$. 
  Der ELBA liefert die Darstellung 
  $$ \ggt (f(x),g(x)) = 1 = r(x)f(x)+s(x)g(x) $$
  mit $r(x), s(x)\in\koerper{p}[x]$. Wird die Berechnung modulo $f(x)$
  ausgef"uhrt, so ergibt sich direkt
  $$ 1\equiv s(x)g(x) \pmod{f(x)} $$
  und $s(x)$ ist somit das inverse Element von $g(x)$ modulo $f(x)$.
  Die Eigenschaft, da\3 $f(x)$ normiert ist, ist unwesentlich, da wir sonst
  mit dem inversen Element des Leitkoeffizients durchmultiplizieren k"onnen.
  Alle Element au\3er der Restklasse der Null haben bez"uglich der 
  Multiplikation ein Inverses. Der Restklassenring ist also ein K"orper 
  mit $p^n$ Elementen, da soviele Restklassen modulo $f(x)$ existieren. 
  \hfill $\Box$

  \begin{satz}[Verfahren von Kronecker]
    Es sei $K$ ein K"orper und $f(x)\in K[x]$ ein nichtkonstantes Polynom.
    Dann existiert ein Erweiterungsk"orper $K$ ($K<L$), so da\3 $f(x)$ eine
    Nullstelle in $L$ besitzt. Ist $f(x)$ irreduzibel, so kann man 
    $$ L:= K[x]/f(x) $$
    setzen.
  \end{satz}
  \beweis Es sei ohne Einschr"ankung $f(x)$ irreduzibel, ansonsten 
  faktorisieren wir $f(x)$ und ersetzen $f(x)$ durch einen 
  seiner irreduziblen Faktoren. Wir bilden die Komposition
  $$ \phi: K \stackrel{\iota}{\hookrightarrow} K[x] 
             \stackrel{\pi}{\rightarrow} K[x]/f(x) =:L. $$
  Der resultierende Homomorphismus $\phi = \pi \circ \iota$ ist injektiv.
  Wir k"onnen $L$ als Erweiterungsk"orper von $K$ auffassen, indem wir
  $K$ mit seinem Bild $\phi (K)$ unter $\phi: K \rightarrow L$
  identifizieren. Wir setzen $\bar{x} := \pi(x)$. Mit 
  $ f(X) = \sum_{i=0}^{n} c_i X^i$ gilt dann
  \begin{eqnarray*}
    f(\bar{x}) & = & \sum_{i=0}^{n} c_i \bar{x}^i \\
               & = & \sum_{i=0}^{n} \pi(x)^i \\
               & = & \pi\left( \sum_{i=o}^{n} c_i x^i\right) \\
               & = & \pi\left(f(x)\right) \\
               & = & \bar{0},
  \end{eqnarray*}
  d.h. $\bar{x}$ ist eine Nullstelle von $f(x)\in K[x]$ in $L$. 
  \hfill $\Box$

  Bei dem Verfahren von Kronecker ensteht $L$ aus $K$ durch
  {\em Adjunktiion einer Nullstelle $\bar{x}$ von $f(x)$}. Die Nullstelle
  $\bar{x}$ wird sozusagen erzwungen, indem wir ausgehend von dem 
  Polynomring $K[x]$ den Restklassenring $L:=K[x]/f(x)$ bilden. Wir
  k"onnen einen Linerfaktor von $f$ "uber $L$ abspalten und k"onnen das
  Verfahren wiederholen. Nach endlich vielen Schritten gelangt man zu einem
  Erweiterungsk"orper, "uber dem $f$ vollst"andig in Linearfaktoren
  zerf"allt.
  \begin{lemma}[Mehrfache Nullstellen]
    Ein normiertes Polynom $f(x)\in K[x]$ hat genau dann keine mehrfachen
    Nullstellen, wenn $\ggt (f(x),f'(x))=1$.
  \end{lemma}

  \begin{definition}[Zerf"allungsk"orper]
    Ein Zerf"allungsk"orper eines Polynoms $f\in K[x]$ ist ein
    Erweiterungsk"orper $L$ des K"orpers $K$, in dem $f$ vollst"andig
    in Linearfaktoren zerf"allt.
  \end{definition}  
  \begin{satz}[Isomorphie der minimalen Zerf"allungsk"orper]
    Es seien $L_1$ und $L_2$ zwei minimale Zerf"allungsk"orper eines
    Polynoms $f\in K[x]$. Dann sind $L_1$ und $L_2$ isomorph.
  \end{satz}
  \begin{lemma}[Gruppenordung]
    Es sei G eine Gruppe. Dann gilt f"ur jedes Gruppenelement 
    $ g^{\betrag{G}} = 1. $
  \end{lemma}
  \beweis mit dem Satz von Lagrange. \hfill $\Box$

  Die multiplikative Gruppe $\koerper{q}^{*}$ eines endlichen K"orpers 
  $\koerper{q}$ hat die Ordnung $q-1$, da alle Elemente au\3er des
  Nullelements invertierbar sind. Daher erf"ullt jedes
  $\alpha\in\koerper{q}$ die Gleichung $X^{q-1}=1$. Wenn man noch das
  Nullelement ber"ucksichtigt, erf"ullt jedes Element 
  $\alpha\in\koerper{q}$ die Gleichung
  $$ f(x) = x\left( x^{q-1} - 1\right) = x^q - x = 0. $$
  Das Polynom $f(x)$ hat $q$ verschiedne Nullstellen in $\koerper{q}$,
  n"amlich die Elemente von $\koerper{q}$. Das Polynom $f(x)$ zerf"allt
  in Linearfaktoren "uber $\koerper{q}$
  $$ f(x) = x^q-x = \prod_{\alpha\in\koerper{q}} \left( x - \alpha \right). $$
  Damit erhalten wir eine zweite Strukturaussage "uber endliche K"orper.
  Alle K"orper $K$ mit $q=p^n$ Elementen sind isomorph, da sie
  minimale Zerf"allungsk"orper des Polynoms $f(x) = x^q - x$ sind.
  
  Diese Aussagen lassen sich zu folgendem Theorem zusammenfassen.
  \begin{theorem}[Strukur endlicher K"orper]
    F"ur jede Primzahl $p$ und jedes $n\in\N$ existiert ein K"orper 
    $\koerper{p^n}$ mit $p^n$ Elementen. Er ist der minimale 
    Zerf"allungsk"orper des Polynoms $x^{p^n}-x$ und seine Elemente sind die 
    Nullstellen dieses Polynoms. 

    Jeder endliche K"orper ist isomorph zu genau einem K"orper
    $\koerper{p^n}$.
  \end{theorem}
  Der folgende Satz ist f"ur die Defintion des diskreten Logarithmus in 
  endlichen K"orpern wichtig.
  \begin{satz}[Struktur der multipikativen Gruppe $\koerper{q}^*$] 
    Die multiplikative Gruppe eines endlichen K"orpers ist zyktisch. Ein 
    Element $\omega\in\koerper{q}^*$ mit 
    $\left<\omega\right>=\koerper{q}^{*}$ hei\3t {\em primitives Element}.
  \end{satz}
%%
%% Index-Calculus
%%
\subsection{Index-Calculus f"ur $\koerper{p}^*$}

Wir stellen die Elemente von $\koerper{p}^*$ als die Menge der Zahlen
$\{1,2,\ldots, p-1\}$. Die Multiplikation wird modulo $p$ ausgef"uhrt.
Der Erzeuger von $\koerper{p}^*$ sei $\alpha$. Es sei $m$ eine Zahl, die
eine Funktion von $p$ ist. Die Faktorbasis $S$ ist die Menge aller
Primzahlen kleiner $m$. 

Der erste und zweite Schritt werden wie im generischen Fall ausgef"uhrt.

Wenn wir die Wahrscheinlichkeit, da"s eine Zufallszahl kleiner $p$ in
Primfaktoren kleiner $m$ zerf"allt, betrachten, k"onnen wir $m$ optimal
w"ahlen, um die Vorberechnungsphase (erster Schritt) und die Berechnung
eines beliebigen Logarithmus (zweiter Schritt) zu beschleunigen. Das f"uhrt
zu einer heuristischen aber subexponentiellen Implementierung des 
Index Calculus Algorithmus mit einer erwarteten Laufzeit von
$L[p,1/2,c]$. Die Laufzeit ist nicht rigoros, da angenommen wird, da"s
das Gleichungssystem im ersten Schritt vollen Rang hat.

\subsection{Index-Calculus f"ur $\koerper{p^n}$}

Die Elemente des endlichen K"orpers $\koerper{p^n}$ werden als Polynome vom 
Grad kleiner $n$ aus $\koerper{p}[x]$ aufgefa\3t. Die Addition ist
gew"ohnliche Polynomaddition und die Multiplikation wird modulo eines
festgelegten irreduziblen Polynoms $f(x)\in\koerper{p}[x]$ ausgef"uhrt.
Wir benutzen also die Isomorphie $\koerper{p^n} \cong \koerper{p}[x]/(f(x))$.

\begin{definition}
  Ein Polynom $w(x)\in\koerper{p}[x]$ hei\3t {\em glatt} bez"uglich
  der Schranke $b\in\N$, genau dann, wenn es in Polynome aus 
  $\koerper{p}[x]$ vom Grad kleiner $b$ faktorisiert.
\end{definition}

Mit der Polynomdarstellung des K"orpers $\koerper{p^n}$ ist es offensichtlich,
die Faktorbasis $S$ als die Menge aller irreduziblen $b$-glatten Polynome aus 
$\koerper{p}[x]$ zu w"ahlen. Im ersten Schritt werden die Logarithmen aller Elemente
der Faktorbasis bestimmt. Es kann gezeigt werden, da"s f"ur kleine $p$ und eine
geeignet Wahl f"ur $b$, der Index Calculus Algorithmus eine erwartete
Laufzeit von $L[p^n,1/2,c]$ hat.

{\bf Algorithmus: }

Es sei $\alpha(x)$ ein Erzeuger von $\koerper{p^n}^*$. Die Faktorbasis
$S$ wird als die Menge aller irreduziblen Polynome $p(x)\in\koerper{p}[x]$ 
mit $\deg(p(x)) \le b$ gew"ahlt. Die Schranke $b$ geben wir sp"ater an.
\begin{enumerate}
  \item Vorberechnungsphase: es werden die diskreten Logarithmen aller 
    Elemente aus der Faktorbasis $S$ berechnet.
    \begin{enumerate}
      \item W"ahle $a_i\in\{1,\ldots,p^n-1\}$ zuf"allig aus und
        berechne $y_i(x) := \alpha(x)^{a_i}$.  
        Wir k"onnen effizient mit dem ggT-Algoritmus "uberpr"ufen, ob
        die Faktoren $p_i(x)\in S$ $y_i(x)$ teilen. Es werden alle
        maximalen Potenzen von $p_i(x)\in S$ von $y_i(x)$ abgespalten.
        Falls wir danach ein Polynom ungleich $c_i\in\koerper{p}$
        erhalten, wissen wir, da"s $y_i(x)$ nicht $b$-glatt ist.
        Falls $y_i(x)$ $b$-glatt ist, so speichern wir die Faktorisierung
        $$ y_i(x) = c_i \cdot \prod_{p_j(x)\in S} p_j(x)^{e_{ij}} $$ mit
        $c_i\in\koerper{p}$. F"ur die diskreten Logarithmen $a_i$
        gilt dann
        \begin{equation} a_i = \dlog{\alpha(x)}{y_i(x)} =
            \dlog{\alpha(x)}{c_i} + \sum_{p_j(x)\in S} e_{ij} \cdot
            \dlog{\alpha(x)}{p_j(x)} \label{loggl} 
        \end{equation}
        Wir wiederholen dies so lange, bis wir mehr als $\#S$ Gleichungen
        der Form~\ref{loggl} haben. 
      \item 
        L"ose dieses lineare Gleichungssystem mit den unbekannten
        $\dlog{\alpha(x)}{p_j(x)}$ f"ur $p_j(x)\in S$.
    \end{enumerate}

  \item Berechnen eines konkreten Logarithmus von $y(x)=\alpha(x)^{s}$. 

        W"ahle $s$ so lange zuf"allig aus, bis $\alpha(x)^s y(x)$
        $b$-glatt ist. Das Polynom $\alpha(x)^s y(x)$ kann dann wie folgt
        $$ \alpha(x)^s y(x) = c \cdot \prod_{p_j(x)\in S} p_j(x)^{e_j} $$
        mit $c\in\koerper{p}$ faktorisiert werden. Daraus kann
        bereits der diskrete Logarithmus $\dlog{\alpha(x)}{y(x)}$ berechnet
        werden, denn es gilt
        $$ \dlog{\alpha(x)}{y(x)} = \dlog{\alpha(x)}{c} - s +
        \sum_{p_j(x)\in S} e_{j} \cdot \dlog{\alpha(x)}{p_j(x)}. $$ 
        Die diskreten Logarithmen $\dlog{\alpha(x)}{p_j(x)}$ f"ur $p_j(x)\in
        S$ wurden im Schritt 1.(a) bestimmt und $\dlog{\alpha(x)}{c}$
        ist im (deutlich kleineren Grundk"orper) einfach zu berechnen.
\end{enumerate}

Die Schranke $b$ sollte so gew"ahlt werden, da\3 f"ur die Menge $S$
$$ \# S \approx \exp\left( c \cdot \sqrt{\log q\log\log q} \right) $$
gilt, wobei $q=p^n$.
%%
%% Coppersmith-Algorithmus
%%

\section{Coppersmith-Algorithmus f"ur $\koerper{2^n}$}

Der Coppersmith-Algorithmus basiert auf der speziellen Darstellung des
K"orpers $\koerper{2^n} \cong \koerper{2}[x]/(x^n + f_1(x))$ mit
$\grad(f_1(x)) \le \log_2(n)$. Ein solches Polynom $f_1(x)$ existiert
f"ur gen"ugend gro\3e $n$.

%%
%% Algorithmus
%%

\subsubsection*{ 1. Erzeugung der Datenbasis}

Es werden die Logarithmen der Elemente aus der Menge $S$ berechnet,
wobei $S$ die irreduziblen Polynome vom Grad kleiner $s$
enth"alt. Dabei wird $s \approx (n\log^2(n))^{1/3}$ und das primitive
Element $\alpha(x)\in S$ gew"ahlt. Beim Index-Calculus-Algorithmus wird
ein beliebiges Polynom (K"orperelement) faktorisiert. Coppersmith
verwendet dagegen zwei Polynome vom Grad $n^{2/3}$. Da diese einen
kleineren Grad haben, steigt die Wahrscheinlichkeit, da\3 sie "uber
$S$ zerfallen.

Es sei $k\in\N$ so gew"ahlt, da\3 $2^k \le (n/\log_2(n))^{1/3}$ und
setze $h:=\oben{\frac{n}{2^k}}$.  W"ahle $u_1(x)$ und $u_2(x)$
teilerfremd vom Grad kleiner $s$ und berechne
\begin{equation} w_1(x) = u_1(x) \cdot x^h + u_2(x) 
\end{equation}
und
\begin{equation} w_2(x) \equiv w_1(x)^{2^k} \pmod{f(x)}. \label{w2} 
\end{equation}
Es gilt dann 
\begin{eqnarray*} 
  w_2(x) & \equiv & u_1(x)^{2^k} \cdot x^{h \cdot 2^k} 
                                      + u_2(x)^{2^k} \pmod{f(x)} \\
         & \equiv & u_1(x)^{2^k} \cdot x^{h \cdot 2^k - n} \cdot f_1(x) 
                                      + u_2(x)^{2^k} \pmod{f(x)},
\end{eqnarray*}
da $x^n \equiv f_1(x) \pmod{f(x)}$ und das Potenzieren mit $2^k$ in
K"orpern der Charakteristik $2$ ein Homomorphismus ist.  Zerfallen
beide Polynome $w_1(x)$ und $w_2(x)$ in irreduzible Polynome aus $S$,
so erhalten wir, wenn wir den diskreten Logarithmus der
Gleichung~\ref{w2} nehmen, die Beziehung
$$ 
\begin{array}{rcll}
  \dlog{\alpha(x)}{w_2(x)} & \equiv & 2^k \cdot \dlog{\alpha(x)}{w_1(x)} 
                  & \pmod{ 2^n - 1} \\
  \Leftrightarrow \dlog{\alpha(x)}{w_2(x)} - 2^k \cdot \dlog{\alpha(x)}{w_1(x)} 
                  & \equiv & 0 & \pmod{ 2^n - 1}.
\end{array}
$$
und damit eine lineare Gleichung f"ur die Logarithmen der Elemente aus $S$.
      
Der Aufwand f"ur die Aufstellung der Datenbasis ist $O(\exp( (c+o(1))
(n\log^2(n))^{1/3} ))$ mit $c \approx \mbox{1,4}$. Diese Laufzeitabsch"atzung
st"utzt sich auf heuristische Annahmen.


\subsubsection*{2. Berechnung des Logarithmus}

Es sei der Logarithmus des Elements $g(x)$ bez"uglich der Basis
$\alpha(x)$ gesucht. Falls $g(x)$ einen Grad $\deg(g(x))\le s$ hat,
dann kann der Logarithmus aus der Datenbasis entnommen werden.

Es sei nun $\deg(g(x))>s$. Wir gehen zun"achst wie beim
Index-Calculus-Algorithmus vor. Es wird $r\in\{1,\ldots,n-1\}$
zuf"allig ausgew"ahlt und
$$ h(x) = g(x) \cdot \alpha^r(x) $$ 
berechnet. Wir faktorisieren $h(x)$ und "uberpr"ufen, ob es glatt
bez"uglich der Schranke $\tilde{s} := (n^2\log_2(x))^{1/3}$ ist. Falls
ja, dann ergibt sich der Logarithmus von $g(x)$ zu
$$ \dlog{\alpha(x)}{g(x)} = \sum_{i=1}^{k}
   \log_{\alpha(x)}{p_i(x)}-r. 
$$ 
Die Polynome $p_i(x)$ sind vom Grad kleiner oder gleich $\tilde{s}$. Die
Logarithmen $\dlog{\alpha(x)}{p_i(x)}$ der Polynome $p_i(x)$ mit
$\deg{p_i(x)}\le s$ stehen in der Datenbasis. Die noch fehlenden
Logarithmen der Polynome $p_i(x)$ mit $s < \deg(p_i(x)) \le \tilde{s}$
werden noch berechnet. Die Berechnung wird vereinfacht, indem man
sie durch Polynome kleineren Grades ersetzen. Zur Bestimmung des
Logarithmus von $p_i(x)$ w"ahlen wir $k$, so da\3 $2^k \le \sqrt{n/b}$
mit $b\approx (n^2\log n)^{1/3} + \deg(p_i(x))/2$ und setzen
$h:=\oben{n/2^k}$.

Anschlie\3end w\"ahlt man zwei zuf"allige teilerfremde Polynome $u_1(x)$
und $u_2(x)$ vom Grad kleiner gleich $\tilde{s}$ und berechnen die
Polynome
$$ w_1(x) := u_1(x)\cdot x^h + u_2(x) $$
und
$$ w_2(x) := w_1(x)^{2^k}, $$
so da\3 $p_i(x) \mid w_1(x)$.
Zerfallen $\frac{w_1(x)}{p_i(x)}$ und $w_2(x)$ in irreduzible Polynome
kleineren Grades,
$$ w_1(x) = p_i(x) \prod_{j=1}^k s_i(x) $$
und
$$ w_2(x) = \prod_{j=1}^l t_j(x), $$
so erh"alt man
$$
\begin{array}{cccc}
   \ds \dlog{\alpha(x)}{w_2(x)}  
      & \equiv 
      & \ds 2^k\dlog{\alpha(x)}{w_1(x)} \equiv \sum_{j=1}^l \dlog{\alpha(x)}{t_j(x)}
      & \pmod{2^n - 1} \\
   \ds 
      & \equiv 
      & \ds 2^k(\dlog{\alpha(x)}{p_i(x)} + \sum_{j=1}^k \dlog{\alpha(x)}{s_j(x))} 
      & \pmod{2^n - 1} \\
   \ds \Rightarrow  \dlog{\alpha(x)}{p_i(x)} 
      & \equiv 
      & \ds 2^{-k} \sum_{j=1}^l \dlog{\alpha(x)}{t_i(x)} - \sum_{j=1}^k \dlog{\alpha(x)}{s_j(x)} 
      & \pmod{2^n - 1}
\end{array}
$$
Diese Gleichnung reduziert die Berechnung der Logarithmen
$\dlog{\alpha(x)}{p_i(x)}$ auf die Berechnung der
$\dlog{\alpha(x)}{s_j(x)}$ und $\dlog{\alpha(x)}{t_j(x)}$. Diese
k"onnen entweder aus der Datenbasis entnommen werden, oder werden
durch rekursive Anwendung des obigen Schrittes berechnet.

Der Aufwand betr"agt $O(\exp(c+o(1)) (n\log^2 n)^{1/3}))$ mit $c
\approx \mbox{1,1}$.
   


      
    
    
        
        
   
   

   
